\documentclass[10pt,a4paper]{article}
\usepackage[utf8]{inputenc}
\usepackage[T1]{fontenc}
\usepackage[french]{babel}
\usepackage[left=1.3cm,right=1.3cm,top=1.3cm,bottom=1.4cm]{geometry}
\usepackage{xcolor}
\usepackage{titlesec}
\usepackage{enumitem}
\usepackage{hyperref}
\definecolor{navy}{RGB}{11,17,32}
\definecolor{cyanacc}{RGB}{2,132,199}
\hypersetup{colorlinks=true,urlcolor=cyanacc,linkcolor=cyanacc,pdfborder={0 0 0}}
\titleformat{\section}{\large\bfseries\color{cyanacc}}{}{0em}{}[\titlerule]
\titlespacing*{\section}{0pt}{1.4ex}{0.8ex}
\setlist[itemize]{leftmargin=1.3em,itemsep=0.25ex,topsep=0.2ex}
\begin{document}
\begin{center}
{\LARGE\bfseries Tsikomia Scidelo ANDRIASAHY}\\[0.5ex]
{\large\color{cyanacc}\bfseries Élève-Ingénieur d'État en Systèmes Électroniques et Automatiques (GSEA)}\\[0.4ex]
{\small Disponible dès Mi-Janvier 2027 • Stage PFE / R\&D (Mobilité Maroc \& International)}\\[0.5ex]
{\small Tanger, Maroc \quad|\quad \href{mailto:scideloandriasahy@gmail.com}{scideloandriasahy@gmail.com} \quad|\quad +212 777 647 362}\\[0.3ex]
{\small \href{https://scidelo-andriasahy.github.io/}{Portfolio Web} \quad|\quad \href{https://www.linkedin.com/in/scidelo}{LinkedIn} \quad|\quad \href{https://github.com/scidelo-andriasahy}{GitHub (19+ Repos)} \quad|\quad \href{https://www.youtube.com/channel/UCqIzu57g0aMRBuQBv294gXw}{YouTube Engineering}}
\end{center}
\section*{Profil \& Objectif PFE}
\noindent Élève-ingénieur en dernière année du cycle Génie des Systèmes Électroniques et Automatiques (GSEA) à l'ENSA de Tanger, je me spécialise à l'interface des systèmes embarqués (STM32G474, DSP TMS320C6713, ESP32), de l'électronique de puissance (KiCad, MicroWind), de la robotique autonome (ROS 2 Humble, Nav2), de l'instrumentation (LabVIEW) et de l'IA industrielle. Fort de deux expériences R\&amp;D chez Fortico (MPPT triphasé entrelacé 4,4 kW à 97,34 \% de rendement, onduleur hybride 230 VAC et 4 cartes PCB IoT) et Smart Automation Technologies (maintenance prédictive par thermographie, modèles de Markov MRF/HMM/HSMM et CNN/U-Net à 98 \%), j'allie conception matérielle rigoureuse et modélisation algorithmique.\par\vspace{0.5ex}
\noindent\textbf{Disponibilit\'e : }Disponible dès mi-janvier 2027 pour un stage de fin d'études (PFE, 4 à 6 mois) ou pré-thèse R\&D. Mobilité totale au Maroc, en France et à l'international.
\section*{Comp\'etences Techniques (Liens Cliquables)}
\begin{itemize}
\item \textbf{Systèmes Embarqués, MCU/DSP \& PCB :} \href{https://www.st.com/en/microcontrollers-microprocessors/stm32g474re.html}{STM32G474}, \href{https://www.ti.com/product/TMS320C6713B}{DSP TMS320C6713}, \href{https://www.espressif.com/en/products/socs/esp32}{ESP32}, \href{https://www.microchip.com/en-us/product/PIC16F877A}{PIC (16F877A)}, \href{https://www.kicad.org/}{KiCad (PCB)}, \href{https://www.labcenter.com/}{Proteus ISIS}, \href{https://fr.mathworks.com/solutions/power-electronics-control.html}{Électronique de Puissance}, \href{https://www.ti.com/product/UCC21520}{UCC21520}, \href{https://www.simcom.com/product/SIM7600X-H.html}{SIM7600}
\item \textbf{Programmation, Firmware \& IoT :} \href{https://misra.org.uk/}{C (Norme MISRA C)}, \href{https://isocpp.org/}{C++}, \href{https://www.python.org/}{Python}, \href{https://standards.ieee.org/ieee/1076/6769/}{VHDL}, \href{https://www.freertos.org/}{Firmware Embarqué Temps Réel}, \href{https://fr.mathworks.com/help/simulink/sfg/what-is-an-s-function.html}{S-Functions C-MEX}, \href{https://platformio.org/}{PlatformIO}, \href{https://www.espressif.com/en/solutions/low-power-solutions/esp-now}{ESP-NOW}, \href{https://mqtt.org/}{MQTT}
\item \textbf{Automatique, Robotique \& IA :} \href{https://fr.mathworks.com/discovery/pid-control.html}{Régulation PID}, \href{https://www.mathworks.com/discovery/genetic-algorithm.html}{GA-PID}, \href{https://fr.mathworks.com/products/fuzzy-logic.html}{Logique Floue (FLC)}, \href{https://www.mathworks.com/help/fuzzy/neuro-adaptive-learning-and-anfis.html}{ANFIS}, MRF / HMM / HSMM, \href{https://pytorch.org/}{Réseaux de Neurones (ANN / CNN)}, \href{https://arxiv.org/abs/1505.04597}{U-Net}, \href{https://arxiv.org/abs/1608.06993}{DenseNet121}, MPPT P\&O Triphasé
\item \textbf{Outils R\&D \& Simulation :} \href{https://fr.mathworks.com/products/simulink.html}{MATLAB / Simulink}, \href{https://docs.ros.org/en/humble/}{ROS 2 (Humble) \& Gazebo}, \href{https://github.com/SteveMacenski/slam_toolbox}{SLAM Toolbox}, \href{http://ceres-solver.org/}{Ceres Solver}, \href{https://navigation.ros.org/}{Nav2}, \href{https://www.ni.com/fr/shop/labview.html}{LabVIEW (Instrumentation)}, \href{https://www.siemens.com/global/en/products/automation/industry-software/automation-software/tia-portal.html}{Siemens TIA Portal (SCADA)}, \href{https://ubuntu.com/}{Linux (Ubuntu)}, \href{https://www.latex-project.org/}{LaTeX}
\end{itemize}
\section*{Projets Techniques \& R\&D (19 Projets — 1 Projet = 1 D\'ep\^ot GitHub)}
\begin{itemize}
\item \textbf{Robot Mobile Autonome AMR (ROS 2 Humble, SLAM Toolbox \& Nav2)} \textit{(Robotique \& ROS 2)} --- Jumeau numérique URDF/Xacro sous Gazebo, cartographie Graph-SLAM (Ceres Solver) et navigation autonome Nav2 validés avec vidéos de simulation et graphes ROS 2. [\href{https://github.com/scidelo-andriasahy/ros2-amr-slam-nav2}{GitHub} | \href{https://scidelo-andriasahy.github.io/assets/pdfs/Rapport_Projet_AMR_ROS2.pdf}{Rapport PDF} | \href{https://scidelo-andriasahy.github.io/?project=1}{Portfolio}]
\item \textbf{Onduleur Hybride 230 VAC \& Convertisseur Solaire MPPT 4,4 kW Entrelacé 3 Phases (Fortico R\&D)} \textit{(Électronique de Puissance)} --- Dimensionnement, schémas KiCad multi-feuilles et simulation MATLAB/Simulink d'un Buck synchrone triphasé entrelacé 4,4 kW (24 VDC / 180 A, rendement 97,34 \%) et d'un onduleur hybride HF 230 VAC pur sinus. [\href{https://github.com/scidelo-andriasahy/mppt-4kw-hybrid-inverter-fortico}{GitHub} | \href{https://scidelo-andriasahy.github.io/assets/pdfs/Rapport_Simulation_MPPT_4kW_Fortico.pdf}{Rapport PDF} | \href{https://scidelo-andriasahy.github.io/?project=2}{Portfolio}]
\item \textbf{Diagnostic Thermique Industriel par IA : Modèles de Markov (MRF, HMM, HSMM) vs Deep Learning (CNN, U-Net)} \textit{(Automatique \& IA Embarquée)} --- Étude comparative R\&D (Smart Automation Technologies) sur 488 thermogrammes FLIR C5 : classification hybride DenseNet121 + KNN/RF à 98 \%, segmentation MRF Potts/ICM vs U-Net (Dice = 0,9851) et chaînes de Markov absorbantes. [\href{https://github.com/scidelo-andriasahy/PROJET_THERMIQUE_2}{GitHub} | \href{https://scidelo-andriasahy.github.io/assets/pdfs/Rapport_PFA_Diagnostic_Thermique_IA_SAT.pdf}{Rapport PDF} | \href{https://scidelo-andriasahy.github.io/?project=3}{Portfolio}]
\item \textbf{Commande \& Suivi de Trajectoire d'un Bras Manipulateur 2 DDL (C-MEX $\times$50, Logique Floue \& DSP TMS320C6713)} \textit{(Automatique \& IA Embarquée)} --- Modélisation lagrangienne accélérée par C-MEX ($\times$50), étude comparative ANN / GA-PID / Logique Floue (49 règles), validation C++ et déploiement sur DSP Texas Instruments TMS320C6713. [\href{https://github.com/scidelo-andriasahy/robot-2dof-fuzzy-gapid-dsp-tms320c6713}{GitHub} | \href{https://scidelo-andriasahy.github.io/assets/pdfs/ROBOT2DOF_Scidelo_FR.pdf}{Rapport PDF} | \href{https://scidelo-andriasahy.github.io/?project=4}{Portfolio}]
\item \textbf{Écosystème IoT Industriel \& Conception PCB sous KiCad : Télématique 4G/GPS, Couveuse PID, Compteur \& IHM ESP32} \textit{(Systèmes Embarqués)} --- Conception matérielle (schémas KiCad, routage PCB 3D, BOM) et spécifications firmware de 4 produits embarqués chez Fortico : Compteur 74HC595, Couveuse PID DRV8825, Boîtier 4G/GPS/RFID et IHM ESP-NOW/MQTT. [\href{https://github.com/scidelo-andriasahy/esp32-industrial-iot-pcb-kicad}{GitHub} | \href{https://scidelo-andriasahy.github.io/assets/pdfs/Schema_KiCad_Smart_Count_ESP32.pdf}{Rapport PDF} | \href{https://scidelo-andriasahy.github.io/?project=5}{Portfolio}]
\item \textbf{Instrumentation Virtuelle \& Oscilloscope Numérique Temps Réel (LabVIEW)} \textit{(Instrumentation \& Signal)} --- Banc de mesure temps réel sous LabVIEW intégrant un générateur de signaux multi-formes et un oscilloscope numérique avec réglage dynamique Volt/Div, ms/Div et curseurs. [\href{https://github.com/scidelo-andriasahy/labview-virtual-oscilloscope-instrumentation}{GitHub} | \href{https://scidelo-andriasahy.github.io/assets/pdfs/RAPPORT_LABVIEW.pdf}{Rapport PDF} | \href{https://scidelo-andriasahy.github.io/?project=6}{Portfolio}]
\item \textbf{Horloge 24H \& Thermomètre Numérique Optimisé en Mémoire (PIC16F877A)} \textit{(Systèmes Embarqués)} --- Firmware C multitâche sur microcontrôleur PIC16F877A avec optimisation algorithmique bas niveau de la mémoire Flash et RAM et pont de puissance H24. [\href{https://github.com/scidelo-andriasahy/pic16f877a-clock-thermometer-optimized}{GitHub} | \href{https://scidelo-andriasahy.github.io/assets/pdfs/RAPPORT_H24_THERMO.pdf}{Rapport PDF} | \href{https://scidelo-andriasahy.github.io/?project=7}{Portfolio}]
\item \textbf{Synthèse Logique Complète d'une Bascule JK \& Caractérisation CMOS 74HC} \textit{(Systèmes Embarqués)} --- Modélisation par machine de Moore, synthèse logique, analyse des marges de bruit, calcul de fréquence maximale (10,5 MHz) et validation réelle en technologie CMOS 74HC. [\href{https://github.com/scidelo-andriasahy/jk-flipflop-cmos-74hc-synthesis}{GitHub} | \href{https://scidelo-andriasahy.github.io/assets/pdfs/Rapport_Bascule_JK_CMOS.pdf}{Rapport PDF} | \href{https://scidelo-andriasahy.github.io/?project=8}{Portfolio}]
\item \textbf{Modélisation \& Commande d'une Machine Asynchrone Triphasée (Park abc/dq0, Gradateurs \& Onduleur V/f)} \textit{(Électronique de Puissance)} --- Simulation numérique complète (RK4) de la machine asynchrone dans les repères stator, rotor et synchrone, puis pilotage par gradateurs (angle de phase / train d'ondes) et onduleur MLI à loi V/f constante avec analyse THD. [\href{https://github.com/scidelo-andriasahy/induction-motor-control-park-inverter-matlab}{GitHub} | \href{https://scidelo-andriasahy.github.io/assets/pdfs/Rapport_Commande_Machine_Asynchrone.pdf}{Rapport PDF} | \href{https://scidelo-andriasahy.github.io/?project=9}{Portfolio}]
\item \textbf{Auto-encodeurs Profonds (Vanilla AE, Sparse SAE, Denoising DAE \& VAE) pour Détection d'Anomalies} \textit{(Automatique \& IA Embarquée)} --- Étude mathématique et expérimentale des architectures d'auto-encodeurs (AE, SAE, DAE, VAE) : régularisation de l'espace latent, débruitage de variété, génération probabiliste (ELBO) et détection non supervisée d'anomalies. [\href{https://github.com/scidelo-andriasahy/deep-autoencoders-anomaly-detection-vae}{GitHub} | \href{https://scidelo-andriasahy.github.io/assets/pdfs/Rapport_Autoencodeurs_IA.pdf}{Rapport PDF} | \href{https://scidelo-andriasahy.github.io/?project=10}{Portfolio}]
\item \textbf{Markov-Vision : Application Desktop Interactive de Segmentation MRF (Gibbs/Potts) \& Synthèse de Texture} \textit{(Automatique \& IA Embarquée)} --- Application graphique Python multithreadée implémentant les Champs Aléatoires de Markov (modèle de Potts) et l'échantillonneur de Gibbs avec suivi temps réel de la convergence énergétique et synthèse de textures. [\href{https://github.com/scidelo-andriasahy/markov-vision}{GitHub} | \href{https://scidelo-andriasahy.github.io/assets/pdfs/Rapport_Markov_Vision.pdf}{Rapport PDF} | \href{https://scidelo-andriasahy.github.io/?project=11}{Portfolio}]
\item \textbf{Interface de Commutation de Puissance NMOS pour Relais Électromagnétiques --- PCB KiCad \& Layout MicroWind} \textit{(Électronique de Puissance)} --- Bureau d'études complet du niveau carte (optocoupleur, diode de roue libre, PCB KiCad \& Gerbers) au niveau silicium (dimensionnement W/L du transistor NMOS en technologie ENSAT et layout physique MicroWind). [\href{https://github.com/scidelo-andriasahy/nmos-power-relay-driver-kicad-microwind}{GitHub} | \href{https://scidelo-andriasahy.github.io/assets/pdfs/Rapport_Interface_NMOS_Relais.pdf}{Rapport PDF} | \href{https://scidelo-andriasahy.github.io/?project=12}{Portfolio}]
\item \textbf{Dimensionnement \& Layout d'un Amplificateur Analogique CMOS Source Commune 200 MHz (MicroWind)} \textit{(Systèmes Embarqués)} --- Conception d'un amplificateur CMOS source commune à charge active PMOS montée en diode respectant le cahier des charges strict : Vdd = 5 V, P $\leq$ 10 mW, Gain G = 5, excursion $\geq$ 0,2 V et bande passante BW = 200 MHz. [\href{https://github.com/scidelo-andriasahy/cmos-common-source-amplifier-200mhz-microwind}{GitHub} | \href{https://scidelo-andriasahy.github.io/assets/pdfs/Rapport_Amplificateur_CMOS_200MHz.pdf}{Rapport PDF} | \href{https://scidelo-andriasahy.github.io/?project=13}{Portfolio}]
\item \textbf{Commande d'un Robot Mobile Unicycle par Logique Floue sous MATLAB / Simulink} \textit{(Robotique \& ROS 2)} --- Modélisation cinématique non-holonome d'un robot mobile unicycle et synthèse de contrôleurs à Logique Floue (.fis) pour le suivi de trajectoire autonome sous MATLAB/Simulink avec vidéo de simulation. [\href{https://github.com/scidelo-andriasahy/unicycle-mobile-robot-fuzzy-simulink}{GitHub} | \href{https://scidelo-andriasahy.github.io/assets/pdfs/Rapport_Robot_Mobile_Unicycle_Flou.pdf}{Rapport PDF} | \href{https://scidelo-andriasahy.github.io/?project=14}{Portfolio}]
\item \textbf{Cinématique Directe/Inverse (FK/IK) \& Génération de Trajectoire d'un Robot 2R sous Simscape Multibody} \textit{(Robotique \& ROS 2)} --- Modélisation mécanique 3D sous Simscape Multibody, résolution des modèles géométriques direct (MGD/FK) et inverse (MGI/IK) et planification de trajectoires articulaires et cartésiennes pour un robot manipulateur 2R. [\href{https://github.com/scidelo-andriasahy/robot-2r-fk-ik-trajectory-simscape}{GitHub} | \href{https://scidelo-andriasahy.github.io/assets/pdfs/Rapport_Robot_2R_FK_IK_Simscape.pdf}{Rapport PDF} | \href{https://scidelo-andriasahy.github.io/?project=15}{Portfolio}]
\item \textbf{Analyse Spectrale FFT \& Détection Automatique des Défauts de Roulement Moteur en Python} \textit{(Instrumentation \& Signal)} --- Application desktop Python/Tkinter de maintenance prédictive : génération de signaux vibratoires/statoriques bruités (SNR réglable), calcul des fréquences caractéristiques de roulement et détection spectrale FFT automatisée. [\href{https://github.com/scidelo-andriasahy/motor-bearing-fault-detection-python-fft}{GitHub} | \href{https://scidelo-andriasahy.github.io/assets/pdfs/Rapport_Detection_Defauts_Roulement_Python.pdf}{Rapport PDF} | \href{https://scidelo-andriasahy.github.io/?project=16}{Portfolio}]
\item \textbf{Conception \& Optimisation de Systèmes Numériques en VHDL sous Intel Quartus Prime} \textit{(Systèmes Embarqués)} --- Synthèse RTL en VHDL, optimisation matérielle (surface logique, fréquence maximale et latence), conception d'automates synchrones et validation fonctionnelle/temporelle sous Intel Quartus Prime. [\href{https://github.com/scidelo-andriasahy/vhdl-digital-systems-quartus-optimization}{GitHub} | \href{https://scidelo-andriasahy.github.io/assets/pdfs/Rapport_VHDL_Quartus_Performance.pdf}{Rapport PDF} | \href{https://scidelo-andriasahy.github.io/?project=17}{Portfolio}]
\item \textbf{Pilotage Industriel, Planification PERT/Gantt \& Suivi Budgétaire EVM sous MS Project} \textit{(Gestion \& Innovation)} --- Gestion complète d'un projet d'ingénierie sous Microsoft Project : structuration WBS, chemin critique PERT/CPM, nivellement des ressources, analyse de trésorerie (cash-flow) et indicateurs de performance EVM. [\href{https://github.com/scidelo-andriasahy/industrial-project-management-msproject-pert-gantt}{GitHub} | \href{https://scidelo-andriasahy.github.io/assets/pdfs/Rapport_MSProject_Gestion.pdf}{Rapport PDF} | \href{https://scidelo-andriasahy.github.io/?project=18}{Portfolio}]
\item \textbf{Continuum Santé : Conception \& Business Model d'une Plateforme E-Santé de Suivi Connecté} \textit{(Gestion \& Innovation)} --- Conception produit, validation d'hypothèses terrain, Business Model Canvas et plan de financement d'une plateforme connectée assurant la continuité de suivi entre patients chroniques et médecins. [\href{https://github.com/scidelo-andriasahy/continuum-sante-connected-health-platform}{GitHub} | \href{https://scidelo-andriasahy.github.io/assets/pdfs/Rapport_Continuum_Sante.pdf}{Rapport PDF} | \href{https://scidelo-andriasahy.github.io/?project=19}{Portfolio}]
\end{itemize}
\section*{Exp\'eriences Professionnelles \& Engagements}
\begin{itemize}
\item \textbf{Ingénieur Stagiaire R\&D --- Électronique de Puissance \& Embarqué} --- \href{https://sitewebfortico.vercel.app/fr}{Fortico} \textit{(Juil. 2026 - Oct. 2026)} : Dimensionnement et simulation MATLAB/Simulink d'un convertisseur solaire MPPT Buck synchrone triphasé entrelacé de 4,4 kW (24 VDC / 180 A) à 97,34 \% de rendement. Conception schématique KiCad d'un onduleur hybride HF 24 VDC / 400 VDC / 230 VAC pur sinus avec synchronisation réseau JIRAMA et anti-îlotage ROCOF. Architecture Dual-Brain (STM32G474 + ESP32) et conception PCB de 4 cartes IoT (Compteur 74HC595, Couveuse PID DRV8825, Boîtier 4G SIM7600/GPS/RFID, IHM ESP-NOW/MQTT).
\item \textbf{Stagiaire Ingénieur R\&D --- IA Industrielle \& IoT (PFA)} --- \href{https://www.linkedin.com/company/smart-automation-technologies/home/}{Smart Automation Technologies} \textit{(Juil. 2026 - Août 2026)} : Benchmark sur 488 thermogrammes (FLIR C5) de moteurs asynchrones : 5 CNN directs (MobileNetV2 94 \%) et 20 hybrides (DenseNet121 + KNN/RF à 98 \%). Segmentation spatiale par MRF Potts/ICM face à U-Net (Dice = 0,9851) et pronostic de défaillance par chaînes de Markov (HMM, HSMM, absorbantes) sur n\oe{}uds IoT.
\item \textbf{Stage Technique en Électricité \& Supervision SCADA} --- \href{https://www.amendis.ma/}{Amendis (Groupe Veolia)} \textit{(Juil. 2025 - Août 2025)} : Immersion au département Électricité : analyse de l'architecture des réseaux de distribution HTB/HTA (20 kV) vers MT/BT. Supervision temps réel au Bureau Central de Conduite (BCC) via systèmes SCADA (Siemens) et protocoles de téléconduite.
\item \textbf{Animateur d'Ateliers Robotique \& Programmation} --- \href{https://www.facebook.com/steamacademy.ma/}{STEAM Academy} \textit{(2025 - Présent)} : Encadrement d'ateliers pratiques de robotique (Arduino), électronique et algorithmique pour jeunes apprenants.
\item \textbf{Professeur Particulier (Sciences \& Technologie)} --- \href{https://www.superprof.ma/}{Superprof / Indépendant} \textit{(2023 - Présent)} : Accompagnement méthodologique et cours particuliers en mathématiques, physique et sciences de l'ingénieur.
\item \textbf{Rédacteur Technique (Freelance)} --- \href{https://www.alphorm.com/}{Alphorm} \textit{(Août 2024 - Mars 2025)} : Rédaction d'articles techniques spécialisés et vulgarisation technologique pour la plateforme Alphorm.
\item \textbf{Adjoint Chargé Académique} --- \href{https://www.linkedin.com/company/eisclubuae/home/}{EISCLUB (ENSA International Students Club)} \textit{(2023 - 2024)} : Intégration académique des étudiants internationaux à l'ENSA Tanger et organisation d'événements scientifiques.
\item \textbf{Responsable Logistique (2 Mandats)} --- \href{https://tsingymaroc.org/}{TSINGY TANGER (Tsingy Maroc)} \textit{(2022 - 2024)} : Coordination opérationnelle, matériel et planification budgétaire des événements associatifs et inter-universitaires.
\item \textbf{Engagements Associatifs, Robotique \& Culturels} --- \href{https://cesamcentrale.org/category/sections-cesam/tanger/}{EnsaTronics, CESAM Tanger, OJA \& FITUT} \textit{(2018 - Présent)} : Membre actif du club robotique EnsaTronics, de la CESAM Tanger et de l'OJA. Dramaturge et acteur au FITUT (Festival International du Théâtre Universitaire de Tanger).
\end{itemize}
\section*{Formation Acad\'emique}
\begin{itemize}
\item \textbf{Cycle d'Ingénieur d'État --- Génie des Systèmes Électroniques et Automatiques (GSEA)} (2022 - 2027 (En cours --- 3ème année Cycle Ingénieur)) --- \textit{ENSA Tanger --- Université Abdelmalek Essaâdi, Maroc}. Systèmes embarqués temps réel (STM32, DSP TMS320C6713, PIC, FPGA/VHDL), électronique de puissance, automatique avancée, robotique autonome ROS 2, instrumentation LabVIEW et IA embarquée.
\item \textbf{Licence 1 --- Mathématiques et Informatique pour les Sciences Sociales (MISS)} (2021 - 2022) --- \textit{Faculté des Sciences --- Université de Fianarantsoa, Madagascar}. Fondamentaux rigoureux en analyse, algèbre linéaire, probabilités, algorithmique et structures de données.
\item \textbf{Baccalauréat Scientifique --- Série C (Mention Très Bien)} (2020 - 2021) --- \textit{Houssen Memorial School (HMS), Tuléar, Madagascar}. Filière scientifique Mathématiques \& Sciences Physiques, Mention Très Bien.
\end{itemize}
\section*{Certifications, Attestations \& Distinctions}
\begin{itemize}
\item \textbf{International Youth Math Challenge (IYMC 2026) --- Qualifié Pre-Final \& Honneur Spécial LaTeX} --- IYMC (ID : PR-2026-1048A24D40BF) (Octobre 2026) [\href{https://scidelo-andriasahy.github.io/assets/pdfs/Certificat_IYMC_2026.pdf}{PDF}]
\item \textbf{Certification en Modélisation \& Simulation (MATLAB / Simulink)} --- IngénieriePro (2025) [\href{https://scidelo-andriasahy.github.io/assets/pdfs/CertificatMatlab.pdf}{PDF}]
\item \textbf{Formation Certifiante en IA \& Idéation Entrepreneuriale} --- Fondation CIDEAL --- INNO'VERT / ENSA Tanger (2025 - 2026)
\item \textbf{Certificat \& Attestation de Stage R\&D --- Électronique de Puissance \& Embarqué} --- Fortico (Direction Générale) (Juil. -- Oct. 2026) [\href{https://scidelo-andriasahy.github.io/assets/pdfs/Certificat_Stage_Fortico.pdf}{PDF}]
\item \textbf{Attestation de Stage PFA --- Diagnostic Industriel par IA (Markov vs CNN)} --- Smart Automation Technologies (SAT) (Juil. -- Août 2026) [\href{https://scidelo-andriasahy.github.io/assets/pdfs/Attestation_Stage_PFA_SAT.pdf}{PDF}]
\item \textbf{Attestation de Stage Direction Électricité --- Réseaux HTA/BT \& SCADA} --- Amendis Tanger (Groupe Veolia) (Juillet 2025) [\href{https://scidelo-andriasahy.github.io/assets/pdfs/Attestation_Stage_Amendis.pdf}{PDF}]
\item \textbf{Certificats de Reconnaissance --- Responsable Logistique (2 Mandats)} --- TSINGY TANGER / Tsingy Maroc (2022 - 2024) [\href{https://scidelo-andriasahy.github.io/assets/pdfs/CertificatTsingyTanger1ereMandat.pdf}{PDF}]
\item \textbf{Distinction Culturelle --- Festival International du Théâtre Universitaire (FITUT)} --- FITUT --- Tanger (2022) [\href{https://scidelo-andriasahy.github.io/assets/pdfs/CertificatFITUT.pdf}{PDF}]
\end{itemize}
\section*{Références techniques \& Bibliographie de travail}
\begin{itemize}
\item \href{https://www.wiley.com/en-us/Analog+Integrated+Circuit+Design\%2C+2nd+Edition-p-9780470770108}{\textbf{Analog Integrated Circuit Design}} --- T. Chan Carusone, D. Johns \& K. Martin (Wiley) \textit{[Microélectronique CMOS]}
\item \href{https://www.dunod.com/}{\textbf{Automatique --- Systèmes linéaires, non linéaires \& représentation d'état}} --- Yves Granjon (2e éd., Dunod) \textit{[Automatique]}
\item \href{https://www.im2np.fr/}{\textbf{Simulation de la variabilité du transistor MOS}} --- Pascal Lemoigne --- Thèse IM2NP \& STMicroelectronics \textit{[Physique MOS]}
\item \href{https://theses.hal.science/tel-00392473v1}{\textbf{Conception de contrôleurs intelligents par approche génétique}} --- Ouahib Guenounou --- Thèse Univ. Paul Sabatier (HAL) \textit{[GA \& Flou]}
\item \href{https://www.mathworks.com/products/fuzzy-logic.html}{\textbf{Optimisation of a Fuzzy Logic Controller Using Genetic Algorithms}} --- Joseph Foran (M.Eng. Report) \textit{[Logique Floue]}
\item \href{https://www.mathworks.com/discovery/genetic-algorithm.html}{\textbf{On-line PID Controller Tuning using Genetic Algorithms}} --- Ian Griffin (M.Eng. Electronic Systems) \textit{[GA-PID]}
\item \href{https://www.mathworks.com/discovery/pid-control.html}{\textbf{The PID Controller Design Using Genetic Algorithm}} --- Saifudin Bin Mohamed Ibrahim (USQ) \textit{[Commande PID]}
\item \href{https://www.wiley.com/}{\textbf{Power Electronics: Converters, Applications, and Design}} --- Ned Mohan, T. Undeland \& W. Robbins (Wiley) \textit{[Électronique de Puissance]}
\item \href{https://www.microsoft.com/en-us/research/publication/pattern-recognition-machine-learning/}{\textbf{Pattern Recognition and Machine Learning}} --- Christopher M. Bishop (Springer) \textit{[IA \& Modèles de Markov]}
\item \href{https://www.st.com/resource/en/reference_manual/rm0440-stm32g4-series-advanced-armbased-32bit-mcus-stmicroelectronics.pdf}{\textbf{STM32G474 Reference Manual (RM0440 --- HRTIM \& Power Control)}} --- STMicroelectronics Official Documentation \textit{[Systèmes Embarqués]}
\item \href{https://navigation.ros.org/}{\textbf{ROS 2 Humble \& Nav2 Official Engineering Documentation}} --- Open Robotics / Nav2 Stack \textit{[Robotique Autonome]}
\item \href{https://ensat.ac.ma/Portail/}{\textbf{Smart Territory Development \& Innovation Challenge (TDC 2026)}} --- Open Innovation \& Agro-industrie durable \textit{[Innovation]}
\end{itemize}
\end{document}